% =====================================================================
% Pre-Oral Compliance Report
% Page numbers are read from the manuscript's .aux file through the xr
% package, so compile Balbin_Archie_W3.tex first; the numbers here then
% always match the current manuscript.
% =====================================================================
\documentclass[12pt,letterpaper]{article}
\usepackage[top=1in,bottom=1in,left=1.25in,right=1in]{geometry}
\usepackage[T1]{fontenc}
\usepackage{helvet}
\renewcommand{\familydefault}{\sfdefault}
\usepackage{array}
\usepackage{longtable}
\usepackage{fancyhdr}
\usepackage{xr}
\externaldocument{Balbin_Archie_W3}

\newcommand{\sysname}{QueuEx}

\pagestyle{fancy}
\fancyhf{}
\fancyhead[R]{\thepage}
\renewcommand{\headrulewidth}{0pt}
\setlength{\parindent}{0pt}
\renewcommand{\arraystretch}{1.25}

% Column types: suggestion and action are justified text, the page
% column is centred; all three are centred vertically in the row.
\newcolumntype{S}{>{\raggedright\arraybackslash}m{5.55cm}}
\newcolumntype{A}{>{\raggedright\arraybackslash}m{5.55cm}}
\newcolumntype{P}{>{\centering\arraybackslash}m{2.3cm}}
\newcommand{\pg}[1]{\large #1}

\begin{document}

\begin{center}
  {\large\bfseries PRE-ORAL COMPLIANCE REPORT}
\end{center}

\vspace{0.4em}
\noindent{\bfseries TITLE: QUEUEX: A COMPUTER VISION-BASED SMART QUEUE
MANAGEMENT AND PREDICTIVE SERVICE DELIVERY SYSTEM}

\vspace{1em}
\begin{tabular}{@{}p{3.6cm}p{0.6cm}l@{}}
RESEARCHERS   & : & ARCHIE D. BALBIN \\
              &   & JOSHUA ROVIC P. SANOTA \\[0.8em]
ADVISER       & : & FRITZGERALD ENRIC T. IMPERIAL, MIT \\[0.8em]
CHAIRMAN      & : & DR. ANNA LORETTA C. ROMULO \\
PANEL MEMBERS & : & DR. RONEL B. SIMON \\
              &   & RONALD O. BALDEMORO, MIT \\[0.8em]
SECRETARY     & : & LIELA B. GARDOSE, MBA \\
\end{tabular}

\vspace{1.2em}
\small
\begin{longtable}{|S|A|P|}
\hline
\multicolumn{1}{|c|}{\bfseries SUGGESTIONS/} &
\multicolumn{1}{c|}{\bfseries ACTION} &
\multicolumn{1}{c|}{\bfseries PAGE} \\
\multicolumn{1}{|c|}{\bfseries RECOMMENDATIONS} &
\multicolumn{1}{c|}{\bfseries TAKEN} &
\multicolumn{1}{c|}{\bfseries NUMBER} \\
\hline
\endfirsthead
\hline
\multicolumn{1}{|c|}{\bfseries SUGGESTIONS/} &
\multicolumn{1}{c|}{\bfseries ACTION} &
\multicolumn{1}{c|}{\bfseries PAGE} \\
\multicolumn{1}{|c|}{\bfseries RECOMMENDATIONS} &
\multicolumn{1}{c|}{\bfseries TAKEN} &
\multicolumn{1}{c|}{\bfseries NUMBER} \\
\hline
\endhead
\hline
\endfoot

\textbf{New Title:}\newline
A Computer Vision-Based Smart Queue Management and Predictive Service
Delivery System
&
Adopted the new title, preceded by the system name \sysname{}, on the
title page and throughout the manuscript.
&
Title page \\ \hline

\textbf{Preliminaries:}\newline
Preliminaries are required during the Pre-Oral Defense.
&
Included the title page, approval sheet, certification pages,
acknowledgment, abstract, table of contents, list of figures, and list
of tables.
&
Preliminary pages \\ \hline

\textbf{Introduction:}\newline
Indent the first paragraph. Use the funnel approach (global situation
to national situation down to local, where the study is conducted).
&
Indented the first paragraph. The Background of the Study now moves from
queue management as a global service-delivery problem, to Philippine
higher education institutions, to Naga College Foundation, where the
study is conducted.
&
\pg{\pageref{cr:background}} \\ \hline

\textbf{Review of Related Literature:}\newline
30 RRLs per topic/theme.
&
Organized the review into four themes, each with at least 30 studies:
person detection (35), face recognition and open-set identification
(30), queue management and re-identification (32), and predictive
analytics for service demand (32).
&
\pg{\pageref{cr:rrl}--\pageref{cr:synthesis}} \\ \hline

\textbf{Synthesis of the State-of-the-Art:}\newline
Indent the first paragraph.
&
Indented the first paragraph.
&
\pg{\pageref{cr:synthesis}} \\ \hline

\textbf{Gap Bridged by the Study:}\newline
Indent the first paragraph. Should summarize the features of the system
presented in the objective of the study.
&
Indented the first paragraph and summarized the system's features in
the order of the objectives: the baseline survey, face-based
identification and queue-number issuance, wait-time prediction measured
by MAE and MAPE, and user satisfaction.
&
\pg{\pageref{cr:gap}} \\ \hline

\textbf{Theoretical Framework:}\newline
Discuss the main theory and the supporting theories. Use the recent
citations of the theories used. Forecast and predictive theories are
outdated; 2004 and 2011 are already outdated.
&
Presented Queuing Theory as the main theory, with Computer Vision, Deep
Face Recognition, and Forecasting Theory as supporting theories. The
2004 and 2011 citations were replaced by Hyndman and Athanasopoulos
(2021), and recent citations were added for each theory (2021--2023).
&
\pg{\pageref{sec:theoretical}--\pageref{sec:conceptual-framework}} \\ \hline

\textbf{Theoretical Paradigm:}\newline
Make it bigger and readable.
&
Redrawn at full page width, with the main theory above the system and the
three supporting theories below it.
&
\pg{\pageref{fig:theoretical-framework}} \\ \hline

\textbf{Conceptual Paradigm:}\newline
Make it bigger and readable.
&
Redrawn at full page width as an Input--Process--Output--Outcome
diagram with a feedback loop, updated to the current system.
&
\pg{\pageref{fig:conceptual-framework}} \\ \hline

\textbf{Statement of the Problem:}\newline
The four research problems given by the panel.
&
Adopted the four research problems. Problem No.~2 refers to the
face-recognition system rather than a YOLO-based system, following the
panel's recommendation to use an appropriate algorithm for facial
recognition: YOLO detects persons, and face recognition identifies
them.
&
\pg{\pageref{cr:sop}} \\ \hline

\textbf{Scope and Delimitation:}\newline
Order of discussion: 1) introduction, 2) scope, 3) limitation.
&
Presented in the required order: introduction, scope, and limitations.
&
\pg{\pageref{sec:scope}} \\ \hline

\textbf{Definition of Terms:}\newline
Indent an introductory statement before defining the terms used in the
study.
&
Added an indented introductory statement before the terms.
&
\pg{\pageref{cr:definitions}} \\ \hline

\textbf{Population and Locale:}\newline
Rephrase the first paragraph, directly addressing NCF as the locale of
the study. How many will be considered as the population? And who are
they?
&
Rephrased the first paragraph to address NCF directly as the locale, and
identified the population as three groups: the cashiering-office staff,
the administrators and IT personnel who oversee the office, and the
students who transact there. The number in each group is to be
confirmed with NCF administration before deployment.
&
\pg{\pageref{cr:population}} \\ \hline

\textbf{Sample Technique:}\newline
What is the sampling technique used?
&
Stated the techniques used: total enumeration for staff and
administrators, stratified purposive sampling guided by Slovin's formula
for students, and convenience sampling for the baseline survey.
&
\pg{\pageref{cr:sampling}} \\ \hline

\textbf{System Architecture:}\newline
Make it readable.
&
Redrawn at a larger scale as six labelled layers, updated to the current
system, and marked by where each layer runs.
&
\pg{\pageref{fig:system-architecture}} \\ \hline

\textbf{Software Development Process:}\newline
Figure 2.3: do not include unless you personalize or customize the
content based on the study's stages and phases.
&
Removed Figure 2.3.
&
\pg{\pageref{sec:sdp}} \\ \hline

\textbf{Validation and Testing Methods:}\newline
If you have a researcher-made questionnaire, include it in the
discussion of reliability testing.
&
Added reliability testing with Cronbach's alpha for both researcher-made
questionnaires. The baseline questionnaire's results are reported in
Table~\ref{tab:sop1-reliability}.
&
\pg{\pageref{cr:validation}} \\ \hline

\textbf{Results and Discussion:}\newline
Does not align with your SOP: 3.1 SOP No.~1, 3.2 SOP No.~2, 3.3 SOP
No.~3.
&
Restructured Chapter~3 so that each section answers one research
problem in order: 3.1 SOP No.~1, 3.2 SOP No.~2, 3.3 SOP No.~3, and
3.4 SOP No.~4.
&
\pg{\pageref{sec:sop1-results}--\pageref{sec:sop4-results}} \\ \hline

\textbf{System:}\newline
Show something ``unique'' in the CV (Computer Vision).
&
Introduced a two-part identity rule: a face is accepted only if its
match score is at least 0.30 \emph{and} it leads the next candidate by
at least 0.15. In testing, the score alone would have given five
strangers a student's queue number; with the margin, none.
&
\pg{\pageref{tab:margin-caught}} \\ \hline

\textbf{System:}\newline
Reduce the 12 algorithms used in the study.
&
Reduced the algorithms from twelve to seven.
&
\pg{\pageref{subsec:cv-layer}--\pageref{subsec:ticket}} \\ \hline

\textbf{System:}\newline
Use appropriate algorithm for facial recognition.
&
Adopted ArcFace face embeddings through the pretrained InsightFace
model, compared by cosine similarity. On held-out photographs, 97.8\% of
enrolled students were linked, none to a wrong identity, and all 45
non-enrolled probes were refused.
&
\pg{\pageref{subsec:reid}, \pageref{subsec:heldout-accuracy}} \\ \hline

\textbf{System:}\newline
Refrain from modifying the current queue flow.
&
Kept the office's existing kiosk and numbering unchanged. Numbers issued
by \sysname{} come from a separate range starting at 5000, so they never
collide with kiosk numbers.
&
\pg{\pageref{subsec:reid}} \\ \hline

\textbf{System:}\newline
For students only. Facial recognition: need to register to validate the
queue number.
&
Only students who register their face in the mobile application with an
NCF Gbox account, and then ask to be queued, are issued a number by the
system.
&
\pg{\pageref{sec:scope}, \pageref{fig:detection-ticket-pipeline}} \\ \hline

\textbf{System:}\newline
Walk-ins (as is printed queue number).
&
Walk-ins continue to take a printed ticket from the existing kiosk,
separately from \sysname{}.
&
\pg{\pageref{sec:scope}} \\ \hline

\textbf{System:}\newline
Conduct load testing/pressure testing.
&
Conducted a load test with 50 concurrent users for two minutes: 2,536
requests with no failures, a median response of 12~ms, and 95\% of
responses within 31~ms.
&
\pg{\pageref{subsec:load-test}} \\ \hline

\textbf{System:}\newline
Include how you train your model, not just discuss the algorithm used.
&
Added the model development: the person-detection and face-recognition
models are pretrained, and their decision thresholds were calibrated on
17 people and 45 held-out photographs. The training of the wait-time
prediction model is reported under SOP No.~3.
&
\pg{\pageref{subsec:calibration}, \pageref{sec:sop3-results}} \\ \hline

\end{longtable}

\end{document}
